\scanpage{47}
% loc: item 2, above first line
\textcircled{2} \inserted{(Via Tate)} \inserted{Soit $\Pi=\Pi_1(S,\unsure{\xi|\zeta})$, soit $\ell$ un nb premier,} À tout $M$, associons
% loc: opening pair
\[
\bigl\{T_{\mathrm B\deleted{\unsure{-\mathrm H|+\mathrm H}}}(M_\xi),\rho_{\ell,M}^{\ell}\bigr\}
\]
% loc: upper left, beside the cancelled block
% slanted
\margin{\begin{tikzpicture}[baseline=(n.base)]\node[draw,ellipse,align=left,inner sep=3pt] (n) {$\mathbb Q_\ell$-Module galoisien\\-- réseau};\end{tikzpicture}}
% loc: two cancelled lines beneath pair
\deleted{où $T_{\mathrm{B.H}}(M_\xi)$ est la structure de \deleted{Betti-}Hodge associée à $M_\xi$, ($T_{\mathrm{B+H}}(M)=\{T_{\mathrm B}(M_{\deleted{\unsure{\xi}}\xi})+\text{bigraduation de }T_{\mathrm D}(M_y)\otimes_{\mathbb Q}\mathbb C\}$,)} \deleted{et} où
% loc: representation beneath cancellation
\[
\rho_{\ell,M}:\Pi\longrightarrow\operatorname{Aut}_{\mathbb Q_\ell}\bigl(T_{\mathrm B}(M_\xi)\otimes_{\mathbb Q}\mathbb Q_\ell\bigr).
\]

% loc: first functor paragraph
Le foncteur ainsi obtenu est pl. fidèle si $S$ est ess. de type fini sur $\mathbb Q$ \deleted{$\mathbb Q$}\inserted{\deleted{un corps $k$ ($\hookrightarrow\mathbb C$ par $\xi$)}}, en admettant comme de juste la conj. de \deleted{Hodge}\inserted{Tate}, et utilisant \unsure{res.|rés.} de Griffiths

% loc: complex base and topological group
Si $S$ est loc de t.f. sur $\mathbb C$, on peut aussi prendre $\Pi'=\Pi_1(S(\mathbb C),\xi)$, et $\rho_M$ au lieu de $\rho_\ell$
\[
\rho_M:\Pi'=\Pi_1(S(\unsure{\mathbb Q|\mathbb C}),\xi)\longrightarrow\operatorname{Aut}_{\mathbb Q}\bigl(T_{\mathrm D}(M_\xi)\bigr),
\]
% loc: left of replacement by B-H
% slanted
\margin{ici c'est conj. Hodge qu'on utilise !}
à condition de remplacer $T_{\mathrm B}$ par $T_{\mathrm{B-H}}$ (structure de Hodge !).

% loc: affine case
Si $S$ est p.ex. affine, \deleted{et} mais pas néc. de type fini sur un corps, on a un résultat analogue, avec $\Pi$ remplacé par un système projectif de grpes fondamentaux de Schém. de t.f. sur $\mathbb Q$ .....

% loc: item 3 and quadruple
3) À tout $M$, associons
\[
\bigl(T_{\mathrm{DR}}(M),\sigma,T_{\mathrm B}(M_\xi),\varphi\bigr)
\]
% loc: description of item 3 and inserted qualifications
où $T_{\mathrm{DR}}(M)$ est la coh. de D.R. de $M$ sur $S$, considéré comme vectoriel \uline{filtré}\inserted{(il suffit de donner la filtration de $T_{\mathrm{DR}}(M)(\xi)$ !)}, $\sigma$ sa stratification de Gauss-Manin relativement à \deleted{$\mathbb Q$}\inserted{un corps de base fixé \uuline{Alg. clos} $k$ (p.ex. $k=\overline{\mathbb Q}$)}, $T_{\mathrm B}(M_\xi)$ la coh. de Betti,
% loc: lower slanted margin
% slanted
\margin{on suppose $S$ donné sur $k$; si $k=\mathbb Q$, il est inutile de mettre la filtration sur $T_{\mathrm{DR}}(M)$, cf plus bas.}
% loc: last comparison and its qualification
$\varphi$ l'isom. canonique $T_{\mathrm{DR}}(M)(\xi)\simeq T_{\mathrm B}(M_{\unsure{y|\xi}})\unsure{\otimes|\ominus}_{\mathbb Q}\mathbb C$. Si $S$ est de t.f. sur $k$, on trouve ainsi un foncteur pl. fid. (moyennant conj. de Hodge\deleted{\unsure{§|?}}, et utilisant Griffiths).
% No watermark lettering securely recoverable on this image; the reading supplied by C remains an audit question.
% archive sheet number: 46

